% Ported from sections/05_conclusion.md (v2, APPROVED 2026-09-18).
% Final-manuscript Sec.~V. Cites greer1993.
\section{Conclusion}
\label{sec:conclusion}

We studied the wetting of a metal film on MgO(001)---the interface at the heart of MgO-based magnetic
tunnel junctions---for the same Fe host with and without added boron. Using a biased,
surrogate-driven search that targets the two configurations the wetting question is about, we mapped
the landscape into its two growth modes: a flat, well-wetting film and a dewetted island. In both
models the island is the ground state, so the flat film is a distinct, higher-energy basin---by
0.1888~\eVA{} for Fe/MgO and by 0.1493~\eVA{} for Fe-B/MgO. Adding boron lowers the flat--island
separation by 0.040~\eVA, a 21\,\% reduction, at a search-level significance of $p = 0.0444$---the
fraction of random relabellings of the 19 searches that would reproduce a shift this large if boron
had no effect. Boron acts inside the film rather than at the interface, and it does not displace the
island as the ground state: the ordered, three-dimensional configuration remains the more stable of
the two in both models.

Boron insertion therefore lowers the relative energy of the flat, well-wetting configuration, moving
it closer to the island ground state without displacing it. This is the direction the confusion
principle of metallic-glass formation predicts---an added element stabilises the disordered, flat
configuration relative to the ordered one \cite{greer1993}---and it matters for MgO-based tunnel
junctions, whose performance depends on a flat, coherent metal film on the barrier. The metal film's
tendency to dewet into islands at monolayer coverage, and the control of interface flatness that
modern junction fabrication demands, make the energetic balance between these two configurations a
quantity worth knowing.
